% PDF 303--341；原书第 275--313 页。
\chapter{检索增强生成}
\label{ch:rag}
\chaptercredit{7}
\chapterabstract{检索增强生成（retrieval-augmented generation，RAG）是对话式大语言模型的一项主要应用。RAG 系统接收用户查询，像聊天机器人一样给出回答，但回答中的事实细节取自文档、结构化数据表等静态知识库。在 RAG 中，先用小型语言模型生成用户查询的嵌入，再与以同样方式嵌入的文档语料比较，找出语义相似、因而与查询语境相关的文本片段。随后，将原始查询和检索到的文档交给经过指令调优的聊天机器人，由它依据文档回答。本章介绍 RAG 的基本原理，说明工程人员构建 RAG 系统时需要考虑的细节。接着深入讨论可纳入 RAG 工作流程的多种模块化改进，以扩展功能并弥补不足。然后介绍评估 RAG 表现常用的若干指标，分别考察稠密检索的准确程度，以及聊天机器人回答查询的成效。最后用 LlamaIndex 的代码演示如何构建、改进和评估一个 RAG 系统。}

% PDF 303--307 上半页；原书第 275--279 页。
\section{引言}
生成式大语言模型的一项有吸引力的应用，是作为通用知识库，接受自然语言查询并给出事实性的回答。如第~3.2.3~节所述，大语言模型将大量信息记在数十亿参数中。用户写好提示，再加上一点运气，有时便能得到准确的回答。不过，训练语料中的信息可能不准确、不完整或缺失，大语言模型也可能产生幻觉（见第~\ref{sec:ch06-hallucination}~节），这些都会影响回答的准确程度。因此，这些信息在现实中能否采用，很大程度上取决于应用能容忍多少错误。教育聊天机器人、医学诊断和自动化客服等应用，可能无法容忍大语言模型回答中的事实错误。

为减少这种回忆不准确或产生幻觉的问题，研究者开发了检索增强生成。最基本的 RAG 方法，用大语言模型为可靠信息库中的文本生成嵌入，针对给定查询快速查找相关段落，并以对用户有用的形式返回信息。从本质上说，RAG 系统是一个问答聊天机器人，从固定的数据库中获取信息，无需依赖模型在预训练时记住的事实细节。这既能提高回答中信息的可靠程度，也能扩展到原大语言模型预训练数据集之外的文档。

\refcite{ch07-ref-011}{Lewis 等人（2020）}最早提出 RAG。ChatGPT 等高性能聊天机器人普及之后，研究人员和开发人员认识到它们较强的上下文推理能力，开始集中研究怎样实现并优化 RAG 框架的各个功能组件，相关研究和技术迅速增加。本章概述 RAG 的要点，讨论近期文献为扩展系统功能、改善系统表现而提出的若干方法，并介绍评估 RAG 系统表现的办法。最后的教程将用常用的 LlamaIndex 包（\refcite{ch07-ref-012}{Liu, 2022}）构建一个 RAG 系统，并试用一些增强方法。

\section{RAG 基础}
\label{sec:ch07-basics}
基本 RAG 系统执行图~\ref{fig:ch07-basic} 中的四个步骤。

\begin{keypoint}
\begin{enumerate}
\item \textbf{建立索引。}将一组文档切分成文本块，用选定的大语言模型把各块转换为文本嵌入。把嵌入放入向量索引（vector index），以便快速与其他向量作语义相似度比较。
\item \textbf{查询。}用户输入一个可以根据文档内容回答的查询。使用构建文档向量索引时所用的同一个嵌入模型，生成查询嵌入。
\item \textbf{检索。}将转换后的查询与向量索引中各个文本块的嵌入比较，通常使用余弦距离，再按与查询的相似程度为文本块排序。取出得分最高的几个块的原文。理想情况下，相似度最高的文本块含有与查询相关的信息。
\item \textbf{生成。}把排名靠前的文本块作为上下文，连同原始查询一起放入提示模板，再将模板交给基于大语言模型的问答智能体。智能体根据检索到的上下文回答问题，并将输出返回给用户。
\end{enumerate}
\end{keypoint}

\begin{figure}[!htbp]
\centering
% 重绘自原书第 277 页的图 7.1。
\begin{tikzpicture}[
  font=\footnotesize\sffamily,
  stepbox/.style={draw=black!53,rounded corners=3pt,text width=37mm,minimum height=19mm,
    inner sep=3.6pt,align=left},
  title/.style={font=\footnotesize\sffamily\bfseries,text=blue!59!black},
  group/.style={draw=blue!58,densely dashed,rounded corners=6pt,inner sep=3.5pt},
  arrow/.style={draw=black!70,-{Stealth[length=1.5mm]},line width=.62pt}]
\node[stepbox,fill=yellow!22] (index) at (0,0) {
  \textbf{建立索引}\\文档分块，语言模型为各块生成向量并存入索引。};
\node[title,above=2pt of index] (setup) {1）构建阶段};
\node[stepbox,fill=cyan!16] (query) at (0,-3.00) {
  \textbf{查询}\\用户输入查询，用同一个语言模型得到查询向量。};
\node[title,above=2pt of query] (ui) {2）用户交互};
\node[stepbox,fill=green!16] (retr) at (6.17,0) {
  \textbf{检索}\\将查询向量与文档向量比较，返回匹配最接近的文档。};
\node[stepbox,fill=violet!15] (gen) at (6.17,-3.00) {
  \textbf{生成}\\把文档与原查询填入提示模板，交给聊天机器人回答。};
\node[title,above=2pt of retr] (runtime) {3）运行阶段};
\begin{scope}[on background layer]
  \node[group,fill=cyan!5,fit=(index)(setup)] {};
  \node[group,fill=red!6,fit=(query)(ui)] {};
  \node[group,fill=yellow!7,fit=(retr)(gen)(runtime)] {};
\end{scope}
\draw[arrow] (index.east) -- ([yshift=4.4pt]retr.west);
\draw[arrow] (query.east) -- ++(.94,0) |- ([yshift=-4.4pt]retr.west);
\draw[arrow,densely dashed] (query.east) -- (gen.west);
\draw[arrow] (retr.south) -- (gen.north);
\end{tikzpicture}

\caption{RAG 系统的基本流程，包括文档向量化与建立索引、用户查询、检索、生成和输出。系统从语料中找出有用的文档，连同原始查询交给生成器，得到依据这些知识的回答。}
\label{fig:ch07-basic}
\end{figure}

图~\ref{fig:ch07-ownership} 给出 RAG 基本流程的一个例子。我们想回答下面的问题。

\begin{quote}
\texttt{Who owns the content created by OpenAI programs?}

OpenAI 程序生成的内容归谁所有？
\end{quote}

如果直接向 ChatGPT 提问，它会这样回答。

\begin{quote}
\texttt{OpenAI typically retains ownership of the content...}

OpenAI 通常保留对内容的所有权……
\end{quote}

原书示例中，这个回答与所用服务条款不符。接下来，取 OpenAI 的服务条款，对文档分块，构建向量索引。把查询转换成嵌入后检索该索引，可以找到两个专门说明 OpenAI 服务输出所有权的文本块。将这两段文档放入固定模板，在提出问题之前把它们作为上下文。把填好的模板交给 ChatGPT，它会这样回答。

\begin{lstlisting}[nolol=true]
As per the provided context information , users own any output
they rightfully receieve from OpenAI services.
\end{lstlisting}

根据所提供的上下文，用户拥有其从 OpenAI 服务正当获得的输出。

\sourcefigure{fig07_02.pdf}{用 OpenAI 输出归属问题演示 RAG 流程。直接询问 ChatGPT 得到了与示例条款不符的回答。用原书所用版本的 OpenAI 服务条款构建向量索引，再作 RAG 查询，得到了符合该版本条款的回答。}{fig:ch07-ownership}

借助 RAG 系统，ChatGPT 现在给出了符合示例条款的回答。完整版教程见第~\ref{sec:ch07-tutorial}~节。

以训练权重之外的信息为根据构建问答系统，可以实现几件事。第一，可以对原先从未进入训练数据的信息，使用大语言模型进行语义推理。这很重要，因为大语言模型的训练集有截止日期，模型无法仅凭预训练回答近期事件的问题。上例中，服务条款在 ChatGPT 的训练阶段之后更新，生成器因而未能回答关于条款的问题。第二，把相关上下文直接传给生成器，要求它只根据这些信息回答查询。与单纯把大语言模型当作知识库相比，这可以改善回答的准确程度和精确性。许多当时表现先进的大语言模型聊天机器人的训练数据集并不透明。模型知道什么，能否可靠地取出相关信息，常常不清楚。模型还容易产生幻觉，以很有把握的语气报告不知道的事情。上下文推理给出的回答，比依赖模型正确表达预训练所得信息更可靠。

\begin{practicaltip}
上述方法仍有若干难点。前面每一步都由多种参数与实现方式控制，配置选得不好，模型便难以达到理想表现。尤其需要确保向量检索能正确找出相关的文本块，并知道怎样询问生成器才能恰当地取出这些信息。缺少这些改进时，RAG 系统也可能像普通的大语言模型调用一样产生幻觉。因此，RAG 系统也必须学会承认不知道答案。出错的代价可能很高。原书举例，加拿大航空公司的人工智能聊天机器人向客户说明了一项不存在的政策，该公司后来被要求履行这一由模型幻觉产生的政策。%
\footnote{\url{https://arstechnica.com/tech-policy/2024/02/air-canada-must-honor-refund-policy-invented-by-airlines-chatbot/}}
\end{practicaltip}

% PDF 307--308；原书第 279--280 页。
\section{优化 RAG}
\label{sec:ch07-optimizing}
RAG 应用的基本结构并不复杂，但有许多细节需要确定。第~\ref{sec:ch07-basics}~节的四个步骤，每一步都有若干预处理操作、架构选择及超参数设置。恰当地选择这些配置，对改善 RAG 系统表现很重要。本节概述基本步骤所含的不同操作，并讨论构建 RAG 系统时应考虑的参数。表~\ref{tab:ch07-enhancements} 汇总了这些细节。首先，构建向量索引时要考虑以下几项。

\begin{itemize}
\item \textbf{文本预处理。}将原始文档转换成合适的格式。Word 文档、HTML 网页、PDF、EPUB 等格式应当转为纯文本。内嵌的表格和 Markdown 格式应当转换成大语言模型易于处理的输入，见后文。也可清理无关信息，将文本规范化，以提高语义匹配的准确程度。对文档进行的规范化，也要在运行时对输入查询采用。
\item \textbf{文本分块。}将文档切分为大小相近的文本块，再生成嵌入。确定每块大小时，需要权衡。块越小，每个向量能表示的信息粒度越细，却可能过小而无法捕捉相关语境。块越小，产生的文本块总数也越多，检索会变慢。可以按固定字符数或词数分块，也可以按句子、段落、节或文档分块，或选择能提高大语言模型编码速度的词元长度。
\item \textbf{元数据。}也可以给文本块补充有用的元数据，包括文档标题、作者、主题、日期或其他用来区分不同文档的类别。可在检索前依据这些信息筛选相关向量，也可以把元数据纳入各文本块，为向量嵌入提供更多语境。
\item \textbf{嵌入模型。}模型的选择决定 RAG 系统能否有效检索到与查询相关的文本块。这是一项语义文本相似性（semantic textual similarity，STS）任务，应选择适合该任务的模型。%
\footnote{Hugging Face 的排行榜用 MTEB 数据集评估 STS 表现（\refcite{ch07-ref-017}{Muennighoff et al., 2023}），可供 RAG 应用参考。见 \url{https://huggingface.co/spaces/mteb/leaderboard}。}
较大的模型通常产生信息更丰富的嵌入，但参数数量、向量维数和生成嵌入所需时间决定计算开销。模型的上下文窗口大小也限制文本块的最大长度。尽管大部分计算开销来自生成文档嵌入，所选模型也决定了运行时生成查询嵌入的办法。
\item \textbf{索引存储。}可从多种方案中选择，需要考虑检索速度、可扩展性、数据库能否扩充、开源或专有，以及集中式或分布式等方面。Superlinked 制作并维护了一份向量数据库比较表，比较各数据库的功能和表现。%
\footnote{\url{https://superlinked.com/vector-db-comparison}}
\end{itemize}

这些步骤决定如何处理和存储文档。下面考察怎样查询这些数据库，以及如何利用检索到的文档生成答案。

\begin{itemize}
\item \textbf{检索函数。}提示与文本块的相似度通常用余弦距离确定。可调节返回的前 $k$ 个文档的数量。$k$ 较小，生成阶段使用的上下文较短，可能更便于模型理解，却可能遗漏得分略低的文档中的相关信息。$k$ 较大，传给生成阶段的信息更多，也更可能被无关信息冲淡。
\item \textbf{生成架构。}生成阶段要决定采用哪个大语言模型、用什么模板组合查询和上下文，以及在这部分内容之前放入哪些系统指令。适用于生成阶段的模型通常是 ChatGPT、Claude 或 Llama-2 这类经过指令调优的大型聊天机器人。成本与质量需要权衡。原书写作时，按词元计算，GPT-4 API 调用费用约为 GPT-3.5-turbo 的 50 倍，但在生成任务上表现更好。
\item \textbf{上下文组织。}需要把可能彼此有差异的前 $k$ 个文档组合成聊天机器人可读的连贯上下文。可以列出全部文档，也可以由生成模型分别摘要，使之更适合上下文窗口，还可以用生成模型将各文本块整合成一段已知信息。
\end{itemize}

这些是构建 RAG 系统时需要考虑的一部分基本事项。下一节将介绍若干可以加入的增强方法，以扩展功能、改善表现并扩大适用范围。

% PDF 309--315；原书第 281--287 页。
\section{增强 RAG}
\label{sec:ch07-enhancing}
2022 年末，ChatGPT 等聊天机器人显著改善了 RAG 应用生成阶段的输出质量。它们还能用接近人类的方式回答事实性问题，使市场广泛需要能够处理训练数据之外信息的聊天机器人。这些变化推动了新 RAG 应用的发展，也催生了可纳入 RAG 工作流程的新模块和新做法。\refcite{ch07-ref-007}{Gao 等人（2024）}将采用这些新方法构建的系统称为“模块化 RAG”（modular RAG）。本节按 RAG 流程的先后顺序，从索引到查询、检索和生成，讨论其中若干改进，重点是每种办法要解决的具体表现问题。图~\ref{fig:ch07-indexing} 至图~\ref{fig:ch07-generation-augment} 将这些优化与增强方法画在 RAG 流程中的大致位置。

\subsection{数据来源与嵌入}
\begin{keypoint}
改进数据索引最直接的办法，是清理数据，例如将输入文本规范化，去掉 HTML 标签等无关标记，以及选择更合适的分块大小。也有较复杂的方法，可以改善模型表现并扩大知识范围。这里简述怎样把结构化数据表纳入 RAG 索引系统，并讨论微调索引嵌入模型的优点。
\end{keypoint}

\subsubsection{使用结构化数据}
\label{sec:ch07-structured}
许多适合用 RAG 查询的信息来源，并不便于转换成纯文本，包括文档中的数据表、SQL 数据库、知识库和网站。数据结构承载大量事实、数值与比较信息，RAG 应用需要能够正确解读。这里回顾纳入这些信息的现有做法。

\refcite{ch07-ref-010}{Hu 等人（2023）}与 \refcite{ch07-ref-024}{Wang 等人（2023d）}等研究考察了在 RAG 检索阶段加入表格查询的办法。例如，可以给一组文档补充一张含有更多相关信息的 SQL 表。再为 RAG 系统配置路由器（见第~\ref{sec:ch07-routing}~节），判断某项用户查询能否利用表中的信息。如果可以，就将用户查询交给在 SQL 代码上训练过的大语言模型，生成符合目的的 SQL 查询。随后用输出查询表格。最后，将返回的信息连同原始查询和其他检索到的文档交给生成器，形成最终回答。

\sourcefigure{fig07_03.pdf}{用于增强 RAG 索引的办法。预处理、分块和元数据（第~\ref{sec:ch07-optimizing}~节）在生成嵌入之前对语料起作用，并可补充结构化数据（第~\ref{sec:ch07-structured}~节）。微调嵌入（第~\ref{sec:ch07-embedding-ft}~节）和选好索引存储（第~\ref{sec:ch07-optimizing}~节），可以改善检索的准确程度和吞吐量。}{fig:ch07-indexing}

举例来说，\refcite{ch07-ref-010}{Hu 等人（2023）}给出下列输入与输出对，演示 ChatDB 系统的能力。

\begin{lstlisting}[nolol=true]
Question: What was the total revenue for January 2023?
LLM-generated Query:
SELECT SUM(total_price) AS total_revenue
FROM sales
WHERE sale_date >= '2023-01-01' AND sale_date < '2023-02-01';
Database response:
+---------------+
| total_revenue |
+---------------+
|         707.0 |
+---------------+
\end{lstlisting}

系统将“2023 年 1 月的总收入是多少？”这一英语请求转换成精确的 SQL 查询，返回所需信息，供生成答案时使用。

\begin{practicaltip}
有些表格没有便于检索的格式。尤其在读入技术 PDF 等文档时，RAG 系统经常遇到包含有用信息的表格，但怎样转换成适合 RAG 的表示并不明确。如果结构不能保留单元格与其标签之间的关系，这些表格就难以发挥作用。为此，人们提出了多种办法，把 PDF 表格转换成便于检索的格式。LlamaParse\footnote{\url{https://www.llamaindex.ai/blog/introducing-llamacloud-and-llamaparse-af8cedf9006b}} 使用专有算法，将各种形状的表格解析成 Markdown 表示，保留表中数值与行、列标签之间的关系。还可以结合为 Markdown 优化的迭代检索方法，提取原表中的数据关系，供生成使用。
\end{practicaltip}

\subsubsection{微调嵌入}
\label{sec:ch07-embedding-ft}
检索的准确程度取决于嵌入模型是否能充分表达 RAG 文档的关键特征，从而取出相关文档。BGE（\refcite{ch07-ref-025}{Xiao et al., 2023}）和 VoyageAI（\refcite{ch07-ref-021}{Wang et al., 2023a}）等开源嵌入模型系列，在语义文本相似性任务上表现出色，但训练语料覆盖范围较广，遇到专业术语与概念较多的主题时，表现可能下降。可以用领域专用样例微调嵌入模型，缓解这一问题。

LlamaIndex（\refcite{ch07-ref-012}{Liu, 2022}）实现了一种常用方法，直接从 RAG 文档构造训练样例。将留出集的文本块交给 GPT-4，让它编写可由对应文档回答的问题。然后微调嵌入，使检索器为每个生成的问题选中正确的来源文档。这样可以使嵌入模型接触专业术语，并弥合查询与 RAG 所选分块形式之间的语义差距。调优之后，可重新生成文档嵌入，再构建 RAG 应用。与使用基础嵌入相比，这一方法使检索准确率提高了 5\% 至 10\%，也改善了在某些专门主题上的表现。%
\footnote{\url{https://blog.llamaindex.ai/fine-tuning-embeddings-for-rag-with-synthetic-data-e534409a3971}}

\subsection{查询}
\begin{keypoint}
RAG 系统的一项主要挑战是，根据人编写的查询找到相关文档。用户之间的措辞差异很大，查询与用来回答查询的文档，在语法和包含的信息上也存在基本差别。这些差异使检索更难。查询增强通常把用户查询变成更容易匹配向量索引中恰当文本块的形式。下面介绍一些解析、转换和扩展查询的办法，以更好地满足用户需求。
\end{keypoint}

\subsubsection{查询改写}
\label{sec:ch07-rewrite}
如本书前面所述，大语言模型的回答可能对输入提示的细节很敏感。措辞、语法等细微变化，交给生成式模型后也可能引起结果的明显变化。RAG 系统需要忠实地回答用户查询，即便查询的写法不够理想。提高成功可能性的一种方法是自动查询改写（query rewriting），将人编写的查询转换成更可能检索到所需结果的提示。

一般无法事先知道哪种提示形式最有效，因此可以按 RAG 流程的结果，用强化学习训练一个改写器来执行这种转换。\refcite{ch07-ref-014}{Ma 等人（2023a）}用 T5（\refcite{ch07-ref-018}{Raffel et al., 2020}）构造查询改写器，再用多个问答训练集与依据生成器输出准确程度构建的奖励模型调优。所得模块位于查询与检索之间，在生成嵌入前，将人编写的查询转换成优化后的形式。实验表明，可训练的改写器优于固定的改写器（定义后未经微调），也优于直接使用原查询。

\sourcefigure{fig07_04.pdf}{RAG 查询阶段的增强方法。查询改写（第~\ref{sec:ch07-rewrite}~节）和查询到文档扩展（第~\ref{sec:ch07-query-to-doc}~节）用调优后的大语言模型改变用户提示，提高检索到正确文档的可能。子查询生成（第~\ref{sec:ch07-subqueries}~节）用模型将复杂提示分解成若干问题，使它们更易在 RAG 数据库中查询。路由（第~\ref{sec:ch07-routing}~节）按查询内容决定采用哪些方法。}{fig:ch07-query-augment}

Raudashcl 提出的 RAG-Fusion 是另一种查询改写方法。%
\footnote{\url{https://towardsdatascience.com/forget-rag-the-future-is-rag-fusion-1147298d8ad1}}
它把数据库的初始查询送给 ChatGPT，改写成若干变体。用各变体分别查询数据库，再通过倒数排名融合（reciprocal rank fusion，RRF）将各次检索所得文档合并为一份排序表。在 RRF 中，某次检索返回的每份文档获得如下分数。
\begin{equation}
\mathrm{Score}_{\mathrm{RRF}}=\frac{1}{r+60},
\label{eq:ch07-rrf}
\end{equation}
其中，$r$ 是该文档在检索结果中的名次。若一份文档在多次检索中出现，就把它的分数加起来。所得分数用于合并成一份排序表。因此，在 RRF 合并结果中得分最高的文档，往往在多个查询变体中都排得较靠前。这样可以减轻选词造成的检索结果波动，返回能够回应多种查询表达方式的文档。

\subsubsection{查询到文档扩展}
\label{sec:ch07-query-to-doc}
基本 RAG 使用查询嵌入搜索一组向量化文本块，取余弦相似度最高的结果，希望它们包含可回答查询的具体信息。一个干扰因素是，查询的语法结构通常与文档文本块不同。如果文本块主题与查询内容足够相似，仍有希望得到合适的匹配，但不同的文本结构可能降低检索效果。

查询到文档扩展（query-to-document expansion）旨在解决这个问题。它先将用户查询交给生成式模型，要求模型构造一段假想文本，其中包含查询的答案。再用嵌入模型将这个文本块向量化，在向量索引中搜索语义相近的文本。这一过程被戏称为生成增强检索（generation-augmented retrieval，GAR）（\refcite{ch07-ref-016}{Mao et al., 2021}）。由于大语言模型可能对查询的答案产生幻觉，生成的文本块常含有错误信息。不过，它仍能形成一段关于查询主题、格式更接近待搜索文档的文本。可以微调生成模型，使输出更像 RAG 文本块，也可以采用少样本上下文学习，将查询与若干示例文本块一起交给模型，让它掌握目标文档的主要特征。\refcite{ch07-ref-006}{Gao 等人（2023）}提出了 HyDE（Hypothetical Document Encoding，假想文档编码），\refcite{ch07-ref-022}{Wang 等人（2023b）}提出了 query2doc，都是这种方法的变体。后者在多种稠密检索任务上，观察到表现最多改善 15\%。

\subsubsection{生成子查询}
\label{sec:ch07-subqueries}
基本 RAG 的一个局限是，只在单条查询上作语义相似度匹配。对于较复杂的查询，回答往往需要综合可能分散于 RAG 语料各处的信息，这个办法常无法找到全部信息。

例如，假设有一组关于美国消费者价格指数（Consumer Price Index，CPI）的报告，用户想知道 2020 至 2023 年各年 3 月的通胀主要由哪些因素推动。这些信息合起来都在文档里，却未必出现在同一个文本块及其周围。需要从各月的 CPI 报告中分别提取。基本 RAG 的语义匹配能力不足以取出所有必要信息，并在生成阶段把它们综合起来。

为弥补这一不足，可以生成子查询（subquery）。使用能力较强的聊天机器人（例如 ChatGPT）把原始查询拆成若干提示，每条提示只针对一项必要信息。每项查询各自执行检索与生成，再把所有回答汇总，作为原始查询的上下文交给最后的生成阶段。对上述 CPI 请求，子查询引擎可能作如下处理。

\begin{lstlisting}[nolol=true]
Input:
What were the key drivers of inflation in the month of March
from 2020-2023?
Subqueries:
1) What were the key drivers of inflation in March 2020?
2) What were the key drivers of inflation in March 2021?
3) What were the key drivers of inflation in March 2022?
4) What were the key drivers of inflation in March 2023?
\end{lstlisting}

这些子查询分别询问 2020、2021、2022 和 2023 年 3 月通胀的主要驱动因素，更能针对性地从 CPI 文档中找到所需信息，提取四个月中各自较重要的通胀成因。再将四个回答合成原始查询的上下文。

\subsubsection{路由}
\label{sec:ch07-routing}
从用户查询到执行检索，RAG 系统可以采用多条途径，例如查询改写、生成子查询，以及查询向量索引以外的知识库。为了利用这些能力，可让 RAG 系统根据查询内容在多种选项之间选择。可以引入路由系统（routing system），接收查询，判断哪些操作最能为好的回答准备信息，再启动相应模块。

通常，这类决策由能力较强的生成式模型完成。为它设计提示模板，让模型仔细考察查询，再从列出的选项中选择。需要考虑的问题包括，查询是否含糊到需要改写，是否必须用多个子提示取回全部必要信息，以及查询是否涉及相关数据库中的信息。

% PDF 315--321；原书第 287--293 页。
\subsection{检索与生成}
\begin{keypoint}
检索的主要目标是为生成器提供回答查询所需的上下文。这个目标有一个重要前提，即按嵌入模型衡量，与查询语义最相近的文本块包含所需信息。如果这个前提不成立，本次 RAG 调用就会失败。检索增强旨在提高所选文档能够回应用户查询的可能性，是改进 RAG 系统较为有效的途径之一。
\end{keypoint}

\subsubsection{重排序}
\label{sec:ch07-reranking}
基本 RAG 的一个常见问题是，最切合查询的文本块往往没有排在语义相似度排序的前列。这种检索误差部分源于常用 RAG 嵌入模型的规模较小。可以用更大、表示能力更强的嵌入模型为语料生成嵌入，改善表现。但许多 RAG 语料规模很大，这样做的费用会很高。

另一个相关问题有时称为“遗失在中间”（lost in the middle，\refcite{ch07-ref-013}{Liu et al., 2023}）。大语言模型较容易准确理解位于提示开头或末尾的上下文信息，较容易遗漏提示中间的信息。如果没有这个问题，只要增加返回的文档数，就有望在结果的某处取得相关信息。但由于中间位置的信息可能被遗漏，这种办法的效果会下降。

重排序（reranking）是在这几项因素之间折中的办法。先用较小的嵌入模型完成初次检索，返回排名靠前的一批文档，例如 20 至 30 份，而基本 RAG 通常只返回几份。接着，用规模大得多、语义表示更丰富的模型重新为这些文档和原始查询生成嵌入，并根据新的向量重新排列前 $k$ 个文本块。

\sourcefigure{fig07_05.pdf}{RAG 检索阶段的增强方法。检索函数（第~\ref{sec:ch07-optimizing}~节）决定收集排名最前的多少份文档。重排序（第~\ref{sec:ch07-reranking}~节）使用第二个、规模更大的嵌入模型，重新排列检索所得文档，使最相关的信息排在前面。迭代检索（第~\ref{sec:ch07-iterative}~节）依次执行多次查询，利用每次返回的文档回答多跳问题。}{fig:ch07-retrieval-augment}

这样，先用成本较低的模型扩大候选范围，再用更好的模型对结果精细排序，便能更准确地选出排在前面的几份文档。这不仅使生成器能用到最相关的文档，也使它们位于列表最前端，避免中间位置的信息被遗漏。较昂贵的模型只用于已经返回的文档，因此可以大幅减轻额外开销，同时使相关文档排在前列。用嵌入对结果重排序并非新办法，\refcite{ch07-ref-015}{Ma 等人（2023b）}等作者推动了这一做法在稠密检索中的应用。LlamaIndex、LangChain 等 RAG 开发软件也广泛实现了它，见第~\ref{sec:ch07-tutorial}~节的教程。

\subsubsection{迭代检索}
\label{sec:ch07-iterative}
有些问题必须综合多项信息才能回答，容易使 RAG 查询失败。需要回答某个主体的事实时，问题可能没有明确提到该主体，而是通过另一个关系间接指向它。\refcite{ch07-ref-020}{Shao 等人（2023）}给出下面的例子。

\begin{quote}
\texttt{Can you get Raclette in YMCA headquarters city?}

在基督教青年会（YMCA）总部所在的城市能吃到 Raclette（拉克莱特奶酪料理）吗？
\end{quote}

在这个例子中，维基百科一类的信息库，可以提供 YMCA 总部在巴黎这一信息，也可以提供 Raclette 是一道在法国同样流行的瑞士菜这一信息。但要碰巧找到一段明确把 YMCA 总部与能否吃到 Raclette 联系起来的文本，并不容易。如果给基本 RAG 生成器一份把 YMCA 与法国关联起来的文档，再给一份把法国与 Raclette 关联起来的文档，它就能回答这个问题。但提示没有明确提到法国，检索时可能漏掉第二项信息。由于必须经过多个推理步骤才能答对，这类问题称为“多跳问题”（multi-hop question）。

\begin{keypoint}
迭代检索（iterative retrieval）是用稠密检索回答多跳推理查询的方法，步骤如下。
\begin{enumerate}
\item 针对查询执行一次检索，返回得分最高的文档。
\item 将查询和文档交给生成器，生成问题的回答。
\item 将回答与原查询组合，用建立索引时所用的语言模型生成嵌入，再查询向量索引（生成增强检索），返回得分最高的文档。
\item 把第 1 步和第 3 步返回的文档合并，连同原查询交给生成器。
\item 请生成模型判断文档及回答是否满足原查询。如果不满足，继续交替进行检索和生成。
\end{enumerate}
\end{keypoint}

\begin{practicaltip}
第二次查询向量索引时，上下文可能包含由幻觉造成的不准确内容。但由于它直接回应查询，有望在语义上接近 RAG 语料中含有正确信息的相关文档。通过多次查询向量索引，迭代检索可以逐步回答问题的不同部分。\refcite{ch07-ref-020}{Shao 等人（2023）}与 \refcite{ch07-ref-005}{Feng 等人（2023）}分别提出了这一思路的实现方法。
\end{practicaltip}

对于上面的例子，\refcite{ch07-ref-020}{Shao 等人（2023）}的迭代检索系统给出了以下结果。

\begin{lstlisting}[nolol=true]
Question: Can you get Raclette in YMCA headquarters city?
*** Iteration 1 ***
-- Retrieved Knowledge --
(1) Title: YMCA Building (Racine , Wisconsin) Context: ... S-
ince the ’50s, headquarters of the publishing house and the
bookstore are located in 11 rue de la Montagne Sainte -Genev-
iève 75005 Paris ...
-- Generation --
The YMCA headquarters is located in Paris, France. Raclette
is a traditional Swiss dish, made with melted cheese served
with potatoes , pickles , and other accompaniments. It is not
a trad- itional French dish, so it is unlikely to be found
in Paris. So the answer is No
\end{lstlisting}

\begin{quote}\small
\textbf{第 1 轮检索}\quad 标题为“基督教青年会大楼（威斯康星州拉辛）”。上下文说，自 20 世纪 50 年代起，出版社总部和书店位于巴黎 75005 区圣热纳维耶芙山街 11 号。

\textbf{第 1 轮生成}\quad YMCA 总部位于法国巴黎。Raclette 是传统的瑞士菜，将融化的奶酪与土豆、腌菜和其他配料一起食用。它不是法国传统菜，所以不太可能在巴黎吃到。因此答案是不能。
\end{quote}

生成器尝试回答巴黎能否吃到 Raclette，却因幻觉得出错误结论。把这次生成的回答与原查询一起转成嵌入，再做一轮检索。

\begin{lstlisting}[nolol=true]
*** Iteration 2 ***
-- Retrieved Knowledge --
(1) Title: Raclette Context: Raclette is a Swiss dish, also
popular in France , based on heating cheese and scraping off
the melted part ...
(2) Title: YMCA Building (Racine , Wisconsin) Context: ... S-
ince the ’50s, headquarters of the publishing house and the
bookstore are located in 11 rue de la Montagne Sainte -Genev-
iève 75005 Paris ...
-- Generation --
YMCA headquarters is located in Paris, France. Raclette is
a dish native to parts of Switzerland , but it is also popu-
lar in France. So it is likely that Raclette can be found
in Paris. So the answer is Yes
\end{lstlisting}

\begin{quote}\small
\textbf{第 2 轮检索}\quad 第一份文档的标题为“Raclette”。上下文说，Raclette 是一道在法国也很流行的瑞士菜，做法是加热奶酪并刮下融化的部分。第二份文档与第 1 轮相同。

\textbf{第 2 轮生成}\quad YMCA 总部位于法国巴黎。Raclette 源自瑞士部分地区，在法国也很受欢迎。因此，很可能在巴黎吃到 Raclette，答案是能。
\end{quote}

第一轮 RAG 调用只返回 YMCA 位置的信息。为了回答查询，生成器对 Raclette 在法国是否流行作了一个没有根据的判断。第二轮把第一轮的生成结果与查询一同嵌入并用于检索。因为这次输入谈到在法国能否吃到这道菜，排名靠前的返回文档也涉及它在当地的流行程度。最后的生成利用了两轮取得的信息，依据上下文得出正确回答。

\subsubsection{上下文整合}
\label{sec:ch07-consolidation}
选好文档后，需要把它们加入模板，再交给生成器。最简单的办法是把各个文本块和提示拼接在一起，由大语言模型整理细节。不过，如果合起来的文本块超过模型的上下文窗口，这个办法就无法使用。如果关键信息所处的位置不合适，模型可能漏掉它，这就是前面所说的中间位置问题。若文本块彼此割裂，还可能缺少把各块信息联系起来的内容。

有多种方法可以更好地综合排名靠前的返回文档中的信息。这一过程称为上下文整合（context consolidation）。常见做法是调用大语言模型，提炼每个文本块的主要事实，缩短交给生成器的上下文（如 \refcite{ch07-ref-003}{Chen et al., 2023b}）。也可以提示模型逐块阅读，每读一块就更新同一份摘要，为全部返回文档构造一份总体摘要（如 \refcite{ch07-ref-026}{Xu et al., 2023}），或采用 LlamaIndex 实现的树形摘要方法\footnote{\url{https://docs.llamaindex.ai/en/latest/examples/response_synthesizers/tree_summarize.html}}（如 \refcite{ch07-ref-012}{Liu, 2022}）。

把一系列相互割裂的检索片段整合成连贯且前后一致的文档，可以改善结果。在多项自然语言处理任务中，\refcite{ch07-ref-026}{Xu 等人（2023）}发现，与直接在提示上下文中拼接返回文档相比，用摘要压缩提示可降低平均困惑度，也就是提高回答的准确程度，同时大幅缩短输入上下文。有些情况下，输入长度可降至原来的 6\%。

\sourcefigure{fig07_06.pdf}{RAG 生成阶段的增强方法。上下文整合（第~\ref{sec:ch07-consolidation}~节）在调用生成器之前，将多份检索文档中的信息提炼成一份文档。上下文组织（第~\ref{sec:ch07-optimizing}~节）为 RAG 系统选择适用的生成提示模板。还可以通过选型甚至微调，优化用于生成的大语言模型聊天机器人（第~\ref{sec:ch07-optimizing}~节）。}{fig:ch07-generation-augment}

% 原书第 292 页。
\begin{table}[p]
\centering
\makebox[\linewidth][c]{\rotatebox{90}{%
\begin{minipage}{.93\textheight}
\caption{检索增强生成各阶段的特点。}
\label{tab:ch07-enhancements}
\footnotesize
\setlength{\tabcolsep}{3pt}
\renewcommand{\arraystretch}{1.17}
\begin{tabularx}{\linewidth}{@{}>{\raggedright\arraybackslash}p{.073\linewidth}
 >{\raggedright\arraybackslash}p{.189\linewidth}
 >{\raggedright\arraybackslash}p{.199\linewidth}
 >{\raggedright\arraybackslash}X
 >{\raggedright\arraybackslash}p{.197\linewidth}@{}}
\toprule
阶段 & 说明 & 优化项 & 干扰因素 & 增强办法\\
\midrule
建立索引 &
把文档切分成文本块，生成各块的嵌入，并存入向量索引。 &
文本规范化\par 表格解析\par 分块策略\par 元数据标签\par 嵌入模型选型\par 向量索引选型 &
数据包含在 PDF 表格中。\par 数据包含在结构化数据库中。\par 文档包含嵌入模型不熟悉的技术语言。 &
用 LlamaParse 做预处理。\par 自动生成数据库查询。\par 微调嵌入模型。\\
\midrule
查询 &
用户输入可以用文档信息回答的问题，再由嵌入模型转换问题。 &
改进查询策略 &
措辞的细微变化可导致不同结果。\par 返回文档相关性低。\par 复杂查询需要多项信息。\par 不同查询需要不同策略。 &
查询改写。\par 查询到文档扩展。\par 生成子查询。\par 路由。\\
\midrule
检索 &
用查询嵌入搜索向量索引，返回语义最相似的文档，按余弦相似度判定。 &
返回排名靠前的文档数\par 对排名靠前的文档重排序 &
最相关的文档未排在前列。\par 多跳问题无法根据一份文档回答。 &
重排序。\par 迭代检索。\\
\midrule
生成 &
将查询和检索到的文档交给聊天机器人，根据检索所得信息生成回答。 &
生成模型选型\par 生成模板和系统提示\par 文档的整合方式 &
生成器输出未忠实采用返回的文档。\par 某些文档的信息在上下文中间被遗漏。 &
改变提示模板。\par 上下文整合。\\
\bottomrule
\end{tabularx}
\end{minipage}}}
\end{table}


\subsection{小结}
构建 RAG 系统时，有许多需要考虑的事项（第~\ref{sec:ch07-optimizing}~节），也有多种可改善特定环节的增强方法或模块。表~\ref{tab:ch07-enhancements} 汇总了这些优化与增强方法。RAG 的组件多种多样，没有一种构建方法能适用于所有情况。每项增强针对某种 RAG 局限加以改进，也各自有适用范围。

\begin{keypoint}
某个 RAG 模块能否发挥作用，要结合系统的使用场景来判断，并与模型复杂度增加所带来的代价权衡。增加模块可能需要投入更多实现与测试工作，增加出错的位置，使响应变慢，还可能因调用更多大语言模型智能体而增加计算和费用开销。%
\footnote{译注：原书将“延迟降低”列为增加模块的代价。根据本段语意，此处按响应变慢处理。}

可以先选取有可能适用的增强方法，再通过试验分别检验各模块对 RAG 表现的影响。下一节讨论评估表现的办法，并说明为了检验 RAG 系统能否达到使用要求，需要考虑哪些实际事项。
\end{keypoint}

% PDF 321--329；原书第 293--301 页。
\section{评估 RAG 应用}
\label{sec:ch07-evaluation}
有效的评估是构建 RAG 方案的重要部分。较复杂的 RAG 应用可能由许多模块组成，往往有不止一个评估目标，评估方法也通常按具体应用的要求选择。

通常，RAG 评估把现有的质量指标和能力指标用于检索器、生成器或端到端输出。有多个评估目标时，要考虑评估对象，以及评估与优化发生在应用开发的哪个阶段，再选择适当的指标和准则。例如，开发早期单独评估生成器，可能会造成过早优化，因为生成器的表现很大程度上受检索器所提供上下文的质量影响，而检索器当时可能还未优化。因此，安排评估步骤的先后，与选对评估指标和能力同样重要（\refcite{ch07-ref-009}{Hoshi et al., 2023}）。

LlamaIndex 文档给出的实用做法是，先做端到端评估。发现局限、边界情况和失败情形后，再逐项评估并优化造成问题的组件（\refcite{ch07-ref-012}{Liu, 2022}）。下面以已有合适的评估流程为前提，考察 RAG 应用常用的评估维度。

\begin{keypoint}
评估 RAG 应用通常有七个方面，可分为三项质量指标和四项系统能力（\refcite{ch07-ref-007}{Gao et al., 2024}）。
\begin{itemize}
\item \textbf{质量指标}
  \begin{enumerate}
  \item 上下文相关性。
  \item 答案忠实度。
  \item 答案相关性。
  \end{enumerate}
\item \textbf{系统能力}
  \begin{enumerate}
  \item 噪声鲁棒性。
  \item 无答案时拒答。
  \item 信息整合。
  \item 反事实鲁棒性。
  \end{enumerate}
\end{itemize}
\end{keypoint}

接下来的两节定义这些方面，说明在典型 RAG 框架的何处以及如何评估它们，并在适当处介绍可用于相应评估的工具和框架。

\subsection{RAG 质量指标}
\label{sec:ch07-quality-metrics}
本节介绍 RAG 的上下文相关性、答案忠实度和答案相关性。图~\ref{fig:ch07-quality} 汇总三者之间的关系。

\begin{figure}[!htbp]
\centering
% 重绘自原书第 295 页的图 7.7。
\begin{tikzpicture}[
  font=\footnotesize\sffamily,
  box/.style={draw=black!59,inner sep=4pt,minimum width=21mm,align=center,
    font=\footnotesize\sffamily\bfseries},
  metric/.style={fill=yellow!26,inner sep=4pt,align=center},
  wire/.style={draw=black!66,-{Stealth[length=1.6mm]},line width=.61pt}]
\node[box,fill=cyan!19] (query) at (0,2.17) {原始查询};
\node[box,fill=green!18] (docs) at (-3.00,0) {检索文档};
\node[box,fill=violet!16] (resp) at (3.00,0) {生成回答};
\draw[wire] (docs.north east) -- node[midway,fill=white,inner sep=1.4pt] {相关于} (query.west);
\draw[wire] (resp.north west) -- node[midway,fill=white,inner sep=1.4pt] {相关于} (query.east);
\draw[wire] (resp.west) -- node[midway,fill=white,inner sep=1.4pt] {忠实于} (docs.east);
\node[metric] at (-3.40,1.68) {上下文相关性};
\node[metric] at (3.40,1.68) {答案相关性};
\node[metric] at (0,-.91) {答案忠实度};
\end{tikzpicture}

\caption{三项 RAG 质量指标之间的关系。上下文相关性衡量检索文档与原始查询的相关程度，答案相关性衡量生成回答与原始查询的相关程度，答案忠实度衡量生成回答对检索文档的忠实程度。}
\label{fig:ch07-quality}
\end{figure}

\subsubsection{上下文相关性}
\label{sec:ch07-context-relevance}
上下文相关性（context relevance）衡量 RAG 检索器能否返回相关上下文，同时舍弃无关上下文。通常可用多项现有指标衡量。有些指标直接考察所有检索到的上下文，不考虑相关性排名，称为与排序无关的指标。另一些指标考虑上下文的相关性排名，称为考虑排序的指标。

\paragraph*{衡量上下文相关性的常用指标}
召回率（recall）是与排序无关的指标，只关注已经检索到的相关上下文数量，以及检索语料中相关上下文的总数。召回率是前者占后者的比例或百分比。检索设置常规定返回的上下文数量上限，因此通常计算召回率@$K$（recall@$K$），其中 $K$ 为取回的固定上下文数量。

\begin{practicaltip}
如果返回上下文的排名影响较小，例如生成阶段使用较短的上下文，或在下游采用重排序器（第~\ref{sec:ch07-reranking}~节），召回率是适用的上下文相关性指标。但如果生成器的提示上下文较长，容易遗漏位于中间的信息，那么只看检索上下文的召回率可能产生误导（\refcite{ch07-ref-013}{Liu et al., 2023}）。计算上下文召回率需要带标签的数据，通常是查询与相关文档的配对。不过，已有实际做法尝试用能力较强的大语言模型半自动计算召回率。例如，使用下面的提示，可让模型依据查询考察上下文。
\begin{lstlisting}[nolol=true]
is all of the required information to answer {query} available in {retrieved_context}
\end{lstlisting}
\end{practicaltip}

精确率（precision）是另一种与排序无关的指标，衡量已检索文档中与查询相关的文档所占的比例或百分比。例如，检索器返回 100 段上下文，其中只有 60 段与查询相关，精确率便为 0.6，即 60\%。该指标有助于了解检索结果中混入了多少无关内容，使开发人员可以着手提高检索精确率，或在下游处理上下文中的额外噪声。下游处理包括检索后重排序（第~\ref{sec:ch07-reranking}~节）、上下文整合（图~\ref{fig:ch07-generation-augment}）和简单筛选规则。和召回率类似，也可以按固定的 $K$ 计算精确率@$K$（precision@$K$）。

\begin{practicaltip}
以往用带标签的查询与相关文档对评估检索上下文的精确率。与召回率一样，这项工作也越来越多地交给能力较强的大语言模型，可采用如下形式的提示。
\begin{lstlisting}[nolol=true]
how many of the search results below are relevant to the
{query}? {retrieved_context}
\end{lstlisting}
模型需要推断具体上下文及其与查询的关系，因此也要考察遗漏中间信息的问题会在多大程度上影响这项指标。
\end{practicaltip}

平均倒数排名（mean reciprocal rank，MRR）衡量在一组查询上，最相关上下文在已排序检索结果中的平均位置。当 MRR 等于 1 时，对于每个受评估的查询，最相关的上下文都排在第一位。当检索器没有返回相关上下文时，MRR 可以为 0。%
\footnote{译注：原书还将最相关上下文每次均排在最后一位列为 MRR 为 0 的情况。若最后一位的排名为有限的 $r$，其倒数仍为 $1/r>0$。}

\begin{practicaltip}
实践中，MRR 通常位于两个极端之间。RAG 应用只向生成器传送一个检索上下文，即 $K=1$ 时，MRR 尤其适合衡量检索效果。它直接反映检索器能否把最相关的上下文排在第一位。把 MRR 与后文的命中率结合起来，可以进一步判断 MRR 较低是因为取回的相关上下文太少，还是因为虽然取回了相关上下文，却排得靠后。例如，若命中率比 MRR 高，可以考虑是检索器的相关性排序不够好，并可在生成前引入重排序器等方法。
\end{practicaltip}

命中率（hit rate）衡量成功检索到最相关上下文的查询所占的比例。实践中，通常对应用选定的 $K$ 个检索上下文计算该指标。命中率容易解释，与 MRR 等指标结合时，可以帮助找出应用中的检索问题。

归一化折损累积增益（normalized discounted cumulative gain，nDCG）是考虑排序的指标，衡量检索上下文的相关程度，并考虑它在结果中的位置。概念上，nDCG 按检索顺序衡量上下文的有用程度，即增益。实际计算时，从结果开头到末尾累加增益。相关性很高却排在后面的上下文所贡献的增益会被折减，再将结果归一化。

使用 nDCG 并解释它与 RAG 结果的关系时，应考虑其假设。尤其是，指标的“折损”部分假定越相关的文档越应该出现在检索上下文的前面。因此，相关性高却排在后面的上下文，在计算中会受到惩罚。如果生成阶段使用所有检索上下文，这个度量可能未必合适，因为那些相关性高却排得靠后的上下文仍会提供给生成器。

\begin{practicaltip}
nDCG 可以理解为检索上下文的排序与理想排序有多接近。在理想排序中，最切题的上下文位于列表顶部，相关性自上而下逐渐降低。因此，nDCG 反映取回的上下文排得好不好。累积增益较低，可能说明上下文需要重新排序。若低分来自相关上下文太少、无关上下文过多，则可能需要提高检索召回率。nDCG 也可以评估 RAG 中的生成器。此时，对给定查询的一组可能回答，按相关性评估排序。
\end{practicaltip}

\subsubsection{答案忠实度}
如第~\ref{sec:ch06-hallucination}~节所述，幻觉是大语言模型的常见缺点。即使表现最先进的大语言模型也会产生幻觉，使得在诚实性十分重要的生产环境中采用它们变得更复杂（见第~5.1.2~节）。RAG 主要就是在大语言模型上加入新的机制，试图缓解这一问题。不过，和任何以概率机制生成输出的系统一样，在开发以及应用继续演变的过程中，都需要评估它能在多大程度上缓解幻觉。

答案忠实度（answer faithfulness）有时也称有据性（groundedness）。它是一个质量指标，用来衡量生成器能否依据检索到的上下文，尽可能准确地陈述事实。实质上，它衡量生成器的回答能在多大程度上以检索器提供的知识为根据。用于评估答案忠实度的指标可分为两类，即基于词语的指标和基于模型的指标。

\paragraph*{基于词语的指标}
基于词语的答案忠实度指标，衡量回答中的词元与提供给生成器的检索上下文或知识重合或以之为根据的程度。在实际应用中，这类办法已大多被后面讨论的基于模型的指标取代。下面列出常用的两种。

\begin{itemize}
\item \textbf{Knowledge-F1 或 K-F1} 用 $F_1$ 衡量生成答案与提供给生成器的检索上下文中的词元重叠。当生成器输出片段的长度与检索上下文长度相差较大时，该指标有局限。例如，生成器给出简洁的回答，而检索上下文通常比回答长。即使回答词元的精确率高，也已正确满足用户需求，仍可能因召回率较低而被扣分（\refcite{ch07-ref-001}{Adlakha et al., 2023}）。
\item \textbf{Knowledge-precision 或 K-precision} 衡量生成器回答中也出现在检索上下文里的词元所占比例。这个指标可以衡量答案忠实度，并避免 K-F1 因度量不对称而可能造成的偏差。此处的不对称在于，忠实回答的事实性词元只能与检索上下文中的事实相同，或只涵盖其中一部分。
\end{itemize}

\paragraph*{基于模型的指标}
在实践中，基于词语的答案忠实度指标已大多被基于模型的方法取代。原因之一是，通过标注生成带标签的上下文较困难，一些指标与人类判断的相关性也较低（\refcite{ch07-ref-001}{Adlakha et al., 2023}）。另一个更重要的原因是，大语言模型处理这类任务的能力不断提高。模型评估还不能解决所有问题（如 \refcite{ch07-ref-023}{Wang et al., 2023c}），但能力最强的大语言模型在判断答案忠实度时，已表现出与人工评估较高的相关性（\refcite{ch07-ref-001}{Adlakha et al., 2023}）。这说明，用较强的大语言模型有望提高 RAG 答案忠实度的评估效率。

Q2 是较早的基于模型的答案忠实度评估方法之一（\refcite{ch07-ref-008}{Honovich et al., 2021}）。

\begin{keypoint}
先从答案中抽取信息丰富的片段，通常用某种命名实体识别（NER）方法。例如，假设 RAG 系统给出下面的回答。
\begin{quote}
\texttt{"Red wine is acidic"}（红葡萄酒是酸性的）
\end{quote}
使用 NER 系统抽取对应于命名实体和名词短语的信息片段，可得到“Red wine”（红葡萄酒）。然后把这些片段与其来源回答一起交给大型生成式模型，以片段及原回答为条件生成相关问题。在上例中，以“Red wine”为信息片段，结合上述回答，可生成这样的问题。
\begin{quote}
\texttt{"What is acidic?"}（什么是酸性的？）
\end{quote}
接着用微调过的问答模型回答这些生成的问题。
\end{keypoint}

原生成系统的答案忠实度，取决于从原回答抽取的信息片段与生成问题的答案有多相似。如果两者完全匹配，Q2 得 1 分。若不完全匹配，则通过自然语言推断（NLI）确定原回答的信息片段与生成问题答案的相似程度。在这一步，蕴含得 1 分，矛盾得 0 分。对于没有给出答案的问答结果，取词元级 $F_1$ 分数。系统的总体 Q2 分数是所有答案对的平均分（\refcite{ch07-ref-008}{Honovich et al., 2021}）。

较新的基于模型的方法利用能力更强的生成式大语言模型，以更集中地衡量答案忠实度。例如，只用 GPT-4，可以大幅减少为问题生成、命名实体识别和问答各准备一个模型的需要（\refcite{ch07-ref-008}{Honovich et al., 2021}）。总体思路与 Q2 相似，但由于 GPT-4 更能利用自然语言理解能力完成整个任务，不必把流程分成那么多模块。

\refcite{ch07-ref-001}{Adlakha 等人（2023）}提出的 LLMCritic 使用一种简单的提示方法，请 GPT-4 评估 RAG 系统的答案是否只包含在检索上下文中出现、或可从中推断的信息和事实陈述。作者给出的提示模板示例如下。

\begin{lstlisting}[caption={评估答案忠实度的提示示例},label={lst:ch07-comparegpt}]
System prompt: You are CompareGPT , a machine to verify the
groundedness of predictions. Answer with only yes/no.
You are given a question , the corresponding evidence and a
prediction from a model. Compare the "Prediction" and the "
Evidence" to determine whether all the information of the
prediction in present in the evidence or can be inferred from
the evidence. You must answer "no" if there are any specific
details in the prediction that are not mentioned in the
evidence or cannot be inferred from the evidence.
Question: {Question}
Prediction: {Model response}
Evidence: {Reference passage}
CompareGPT response:
\end{lstlisting}

这个提示要求 GPT-4 核查 RAG 答案中的全部信息，判断它们是否已在材料，即检索上下文中出现，或能否从该材料推断出来。如果预测包含材料没有提及或无法推断的具体细节，则回答“no”。一般而言，用大语言模型评估答案忠实度时，指标计算公式为
\begin{equation}
\text{忠实度}=
\frac{|\text{答案中可由检索上下文推断的事实数}|}
     {|\text{答案中的事实总数}|}.
\label{eq:ch07-faithfulness}
\end{equation}

随着大语言模型能力提高，越来越多的这类 RAG 评估指标用大语言模型计算。LlamaIndex（\refcite{ch07-ref-012}{Liu, 2022}）、TruLens（\refcite{ch07-ref-019}{Reini et al., 2024}）和 Ragas（\refcite{ch07-ref-004}{Es et al., 2023}）等 RAG 评估框架或工具提供了多种这类方法，后面的教程将演示其中一些。

\subsubsection{答案相关性}
答案相关性（answer relevance）是评估 RAG 系统质量的另一项重要指标，回答“生成的答案与用户查询和检索上下文有多相关？”最常见的计算方法也使用能力较强的大语言模型。将生成器给出的答案、生成该答案所用的上下文，以及要求根据这些信息生成 $N$ 个合成问题的指令，交给评估模型。然后，在语义上将这些问题与最初的用户查询比较。用于生成合成问题的答案，正是为响应这条原查询而给出的。答案相关性取原始查询与 $N$ 个合成问题之间的平均语义相似度。因而，如果根据 RAG 答案生成的问题在语义上更接近原查询，该答案的相关性分数就更高，反之则更低。具体来说，Ragas 框架用下式计算这个指标。
\begin{equation}
\text{答案相关性}=\frac{1}{N}\sum_{i=1}^{N}
  \operatorname{sim}(E_{g_i},E_o),
\label{eq:ch07-answer-relevance}
\end{equation}
其中，$N$ 是评估模型生成的合成问题数，$E_{g_i}$ 是第 $i$ 个合成问题的嵌入，$E_o$ 是原始查询的嵌入，$\operatorname{sim}$ 是两者之间合适的相似度度量，例如余弦相似度。

答案相关性可以通过与真实参考答案比较等方式评估。不过，对输出形式丰富的应用，定义真实参考答案通常复杂且开销较高。大语言模型评估器可以减少这项工作，因此正在成为常用方法。

% PDF 329--335；原书第 301--307 页。
\subsection{评估 RAG 系统能力}
除了用定量指标直接评估 RAG 系统的检索质量、答案忠实度和相关性，还需要考察几项能力，以了解系统中大语言模型生成器的总体表现、适用范围和可靠程度（\refcite{ch07-ref-002}{Chen et al., 2023a}）。这类评估有助于弄清检索上下文和提示扩充怎样影响模型生成。例如，要了解生成器在多大程度上采用非参数知识，即提示中的上下文信息，替代参数知识，即预训练时存入模型参数的信息。如果生成器用检索到的非参数知识覆盖参数知识，而检索结果有误，开发人员就可能需要修改提示，或提高检索精确率，使生成的回答尽量符合事实。

下面介绍常用的四项 RAG 能力评估。图例取自 \refcite{ch07-ref-002}{Chen 等人（2023a）}，图~\ref{fig:ch07-capabilities} 汇总各项能力。

\subsubsection{噪声鲁棒性}
噪声鲁棒性（noise robustness）衡量大语言模型生成器能否在含有噪声的检索文档中，只采用有用信息，如图~\ref{fig:ch07-noise} 所示。它考察生成器在上下文混有无关信息时，能否给出用户查询的正确答案。

\sourcefigure{fig07_08.pdf}{RAG 回答中的噪声鲁棒性示例。虽然检索上下文含有噪声，例如关于诺贝尔物理学奖而非文学奖的信息，生成器仍能给出用户查询的正确答案。}{fig:ch07-noise}

\begin{practicaltip}
要评估 RAG 大语言模型处理含噪上下文的能力，需要知道对于一组生成的问答对，哪些上下文相关，哪些不相关。通常将相关文档与按适当比例随机抽取的不相关上下文配对。再将这些“检索”上下文交给生成器，并将它的答案与“问题—支持信息—答案”三元组数据集中的原答案比较。可用精确匹配（exact match，EM）等准确性指标量化模型在这方面的能力（\refcite{ch07-ref-002}{Chen et al., 2023a}）。
\end{practicaltip}

\subsubsection{无答案时拒答}
在无答案时拒答（negative rejection）这一能力中，如果没有任何检索上下文包含回答用户查询所需的相关信息，RAG 应用就应拒绝作答。如图~\ref{fig:ch07-rejection} 所示，面对“谁获得了 2022 年诺贝尔文学奖？”这一问题，所有上下文都缺少回答所需的事实，取回的内容仅与诺贝尔物理学奖有关。

评估这种能力，有助于开发人员改进应用在知识来源不足以支持忠实或事实正确的回答时的行为。例如，可以针对这种情况设置更明确的系统指令。%
\footnote{译注：原书图 7.9 的图注写作 2021 年，图内的用户查询和正文都写作 2022 年。}

\sourcefigure{fig07_09.pdf}{没有任何检索上下文包含回答所需信息时，如果生成器不给出用户查询的答案，就表现出无答案时拒答的能力。本例的上下文没有包含 2022 年诺贝尔文学奖得主的信息，因此生成器表示无法回答。}{fig:ch07-rejection}

\begin{practicaltip}
评估无答案时拒答的能力，同样先生成一组问答对及支持答案的信息。然后只从不相关文档抽取上下文，即抽取那些并非随问答对生成的文档。根据模型正确给出特定拒答内容的次数，可衡量 RAG 大语言模型的拒答能力。拒答内容通常由给模型的指令指定。计算拒答率，可了解查询 RAG 系统时正确和错误拒答各占多少。
\end{practicaltip}

\subsubsection{信息整合}
RAG 中的信息整合（information integration）指大语言模型生成器能从多份上下文文档中综合信息，为复杂问题给出正确答案。如果用户查询包含多个子问题，或涉及分布在各处的多项信息，RAG 系统需要将各项信息对应的相关上下文整合进回答。如图~\ref{fig:ch07-integration} 所示，用户询问 2021 年和 2022 年诺贝尔物理学奖的得主。回答所需的外部信息也分布在两份上下文文档中。RAG 回答正确地综合了这些上下文，给出了所需的答案。

\sourcefigure{fig07_10.pdf}{RAG 生成器能利用多份上下文文档回答复杂问题时，表现出信息整合能力。本例中，用户同时询问 2021 年和 2022 年诺贝尔物理学奖的得主，生成器综合相关上下文，给出正确答案。}{fig:ch07-integration}

\begin{practicaltip}
实践中，评估信息整合也依赖于生成“问题—支持信息—答案”三元组。但在构造测试数据时，要多做一步，为问题增加应回答的方面。例如，可以合并两个问题及其各自的答案和支持信息，使回答这个较复杂的问题所需的信息分散在多份上下文文档中。也可用精确匹配等准确性指标判断信息是否成功整合，直接比较 RAG 生成的回答与原先为问题生成的答案（\refcite{ch07-ref-002}{Chen et al., 2023a}）。
\end{practicaltip}

\subsubsection{反事实鲁棒性}
RAG 应用常用的外部知识库也可能含有事实错误。因此，需要评估 RAG 生成器能否识别检索上下文中的错误，这称为反事实鲁棒性（counterfactual robustness）。识别检索上下文中的错误，完全依赖生成器的参数知识。如果所选大语言模型没有相应领域的知识，或掌握得不充分，这项能力就难以评估。可以对大语言模型做领域适配或微调，但费用较高，也会抵消 RAG 的一些优点。不过，开源领域已有许多经过领域微调的大语言模型，因此生成适用的领域测试数据也变得更可行。

\sourcefigure{fig07_11.pdf}{反事实鲁棒性指生成器能检测所提供上下文的事实错误，并在回答中指出错误。这项能力依靠大语言模型生成器的参数知识。如果模型缺少相应领域知识，或应用依赖模型知识截止日期之后的信息，这项能力可能难以评估。}{fig:ch07-counterfactual}

\begin{practicaltip}
检验这项能力时，先提示大语言模型生成器只根据其参数知识生成问题和答案。也就是说，让模型提出它已经知道答案的问题，无需依赖外部知识库。为构造反事实支持信息，先核实生成的答案正确，再从外部知识中检索支持信息。然后手工修改支持上下文，用错误的事实代替正确的支持事实。用这个数据集考察两项指标，即错误检测率和错误纠正率。

错误检测率取决于支持上下文含有事实错误时，生成器是否返回指定内容。它衡量大语言模型能在多大程度上根据参数知识判断检索上下文的事实性。类似地，错误纠正率衡量支持信息包含错误时，生成器仍能给出正确答案的频率。
\end{practicaltip}

\begin{figure}[!htbp]
\centering
% 重绘自原书第 306 页的图 7.12。
\begin{tikzpicture}[
  font=\footnotesize\sffamily,
  topbox/.style={fill=yellow!26,minimum width=24mm,minimum height=8mm,
    inner sep=2pt,align=center},
  band/.style={draw=black!78,minimum width=109mm,minimum height=8.7mm,
    align=center,font=\sffamily\bfseries},
  line/.style={draw=black!71,line width=.56pt},
  arrow/.style={line,-{Stealth[length=1.5mm]}}]
\node[band,fill=violet!17] (response) at (0,0) {生成回答};
\node[band,fill=green!18] (documents) at (0,-3.67) {检索文档};
\foreach \x/\item in {-3.99/噪声鲁棒性,-1.33/无答案时拒答,1.33/信息整合,3.99/反事实鲁棒性}{
  \node[topbox] (metric) at (\x,2.32) {\item};
  \draw[line] (metric.south) -- (response.north -| metric);
}
\node[fill=white,inner sep=2pt,font=\scriptsize\sffamily] at (0,1.15)
  {衡量生成回答的下列能力};
\foreach \x/\y/\item in {
   -3.99/-1.36/忽略其中的\\无关信息,
   -1.33/-2.44/承认其中缺少\\回答所需的信息,
   1.33/-1.36/整合来自\\多份文档的信息,
   3.99/-2.44/识别并拒绝\\其中的错误信息}{
  \coordinate (lane) at (\x,0);
  \draw[arrow] (response.south -| lane) -- (documents.north -| lane);
  \node[fill=white,inner sep=1.8pt,align=center,font=\scriptsize\sffamily] at (\x,\y) {\item};
}
\end{tikzpicture}

\caption{RAG 系统各项能力指标所衡量的性质。四项指标各自衡量生成器能否理解检索文档所具有的正面或负面性质，并作出合适的回答。}
\label{fig:ch07-capabilities}
\end{figure}

\subsection{RAG 评估小结}
RAG 应用有多个边界不容易划清的功能组件，因此评估较为复杂。开发人员需要考虑向量数据库的索引是否合适、外部知识文档如何分块、检索结果的相关性与排序，以及最后的生成阶段。已有许多研究与工具简化了这个过程。从 TruLens 团队提出的 RAG 三项指标框架，到 LlamaIndex 等便于实现的框架，这些工具减少了需要用户自行处理的细节，使其能够快速制作原型，并开发适合生产使用的应用。

第~8~章将讨论在生产环境中使用大语言模型时涉及的运维概念、框架、工具与挑战，其中不少也适用于 RAG 应用开发。深入讨论这些问题之前，先给出一个 RAG 开发和评估的教程。

% PDF 335--339；原书第 307--311 页。
\section{教程——构建自己的检索增强生成系统}
\label{sec:ch07-tutorial}
\subsection{概述}
检索增强生成可以结合大语言模型的语义理解能力和直接来源材料的事实准确性。本章从基本做法、优化到核心功能的增强，再到表现评估，讨论了这类系统的多个方面。本节以一个常用的开源 RAG 应用为例，演示怎样实际构建和增强 RAG 系统。

近几年，出现了多种这类开源库。它们提供高级功能，实现不同的 RAG 方法，并允许在构建自己的系统时充分定制。本教程用 LlamaIndex（\refcite{ch07-ref-012}{Liu, 2022}）构建 RAG 系统，试验部分可调参数，并评估系统表现。

\begin{keypoint}
\noindent\textbf{目标}
\begin{itemize}
\item 演示如何用 LlamaIndex 轻松搭建一个基本的 RAG 应用。
\item 探索定制和改进 RAG 应用的多种可能。
\item 评估 RAG 应用的上下文相关性、答案相关性和有据性。
\end{itemize}
\end{keypoint}

这里给出教程的简明版本。完整版见 \url{https://github.com/springer-llms-deep-dive/llms-deep-dive-tutorials}。

\subsection{实验设计}
这个练习按步骤构建一个实验性 RAG 应用。文档语料使用 2024 年 1 月下旬的 OpenAI 条款与政策，取自 \url{https://openai.com/policies}。应用使用以下工具。

\begin{itemize}
\item \textbf{LlamaIndex。}该框架提供许多模块，易于纳入同一个 RAG 框架，处理文档解析、索引、搜索与生成。它与 Hugging Face 集成，便于灵活选择嵌入和生成模型（\refcite{ch07-ref-012}{Liu, 2022}）。
\item \textbf{\texttt{BAAI/bge-small-en-v1.5}。}北京智源人工智能研究院（Beijing Academy of Artificial Intelligence，BAAI）文本嵌入模型系列中的小型英语版本。它在文本相似性任务上表现良好，参数量约为 3340 万，规模较小，便于微调。
\item \textbf{OpenAI ChatGPT。}整个教程使用 OpenAI 的 \texttt{gpt-3.5-turbo} 与 \texttt{gpt-4} 模型作生成器，也用它们来比较 RAG 系统的输出。
\end{itemize}

首先，把 OpenAI 条款与政策中的各份文档载入 LlamaIndex。接着，选择分块策略和嵌入模型来生成向量索引。完成后，就可以通过 LlamaIndex 查询 RAG 数据库，并用 \texttt{gpt-3.5-turbo} 作生成模型。

以初始应用为起点，接下来探索多种设计选择和改进。其中较重要的两项是微调嵌入模型和加入文档重排序器。每引入一种做法，就比较结果，考察应用有什么变化。最后，将改进后的 RAG 应用与未微调嵌入模型的早期版本作全面比较。

\subsection{结果与分析}
评估方法有很多，这里只考察第~\ref{sec:ch07-quality-metrics}~节的三项质量指标。

\begin{enumerate}
\item \textbf{上下文相关性。}检索到的上下文是否与查询相关？
\item \textbf{答案相关性。}生成的答案是否与查询相关？
\item \textbf{答案忠实度或有据性。}生成的答案是否有检索上下文支持？
\end{enumerate}

用这三项指标检查查询能否找到有用文档、生成的回答是否忠实于文档，以及它是否回答了原查询。

先看上下文相关性。从大量语料块生成问题，构造测试集。再用 RAG 为各个生成的问题检索文档，检查该问题的来源文档是否被选中。这里使用两项指标。一项是 \texttt{mrr}，即平均倒数排名，对单个问题为 1 除以正确文档的名次。另一项是 \texttt{hit\_rate}，若正确文档位于前五项中，取 1，否则取 0（见第~\ref{sec:ch07-context-relevance}~节）。以下是一个问题的结果。

\noindent\begin{minipage}{\linewidth}
\begin{lstlisting}[caption={上下文相关性示例},label={lst:ch07-context-example}]
Total MRR = 1.0 / 1
# Hits = 1.0 / 1
Expected ID =
['f712129a -a58d-4e36-b62f -22ebfeda56a8 ']
Retrieved IDs =
['f712129a -a58d-4e36-b62f -22ebfeda56a8 ',
'1f95b363 -1002-4f2d-bf17-ab4842714072 ']
\end{lstlisting}
\end{minipage}\par

可见，预期的文档 ID 在检索结果中排名第一，因此 MRR 和命中数均为 1/1。再看由 50 个验证问答对组成的样本。

\noindent\begin{minipage}{\linewidth}
\begin{lstlisting}[caption={上下文相关性分数},label={lst:ch07-context-scores}]
- Base model:
  - Total MRR = 36.5 / 50
  - # Hits = 42.0 / 50
- Fine-tuned model:
  - Total MRR = 40.0 / 50
  - # Hits = 46.0 / 50
\end{lstlisting}
\end{minipage}\par

多数情况下都能返回来源文档，而且它经常排在首位，但并非总是如此。此前用 OpenAI 条款与政策语料微调过嵌入模型的 RAG 系统表现略好。

接着考察答案相关性，即 RAG 流程能否为查询给出切题的回答。提交一个查询，收到回答后，请 GPT-4 判断回答是否切合问题。这里返回简单的 \texttt{True} 或 \texttt{False}。以下是一项测试。

\begin{lstlisting}[caption={答案相关性示例},label={lst:ch07-answer-example}]
query = "How can individuals request corrections for factually
inaccurate information about themselves in ChatGPT output?"
results = run_answer_relevance_eval(index, [query ,])
Response:
Individuals can request corrections for factually inaccurate
information about themselves in ChatGPT output by submitting
a correction request through privacy.openai.com or by sending
an email to dsar@openai.com. If the inaccuracy cannot be
corrected due to the technical complexity of the models , they
can request the removal of their Personal Information from
’ChatGPTs output by filling out a specific form.
Relevant:
True
\end{lstlisting}

\begin{quote}\small
\textbf{查询}\quad 如果 ChatGPT 输出了关于某个人的不准确信息，当事人如何请求更正？

\textbf{回答}\quad 当事人可以通过 \nolinkurl{privacy.openai.com} 提交更正请求，或发送邮件至 \nolinkurl{dsar@openai.com}。如果模型的技术复杂性使错误无法更正，当事人可填写相应表格，请求从 ChatGPT 输出中移除其个人信息。

\textbf{是否相关}\quad 是。
\end{quote}

这一回答确实回应了查询。用验证集的 50 个样本评估，结果如下。

\noindent\begin{minipage}{\linewidth}
\begin{lstlisting}[caption={答案相关性分数},label={lst:ch07-answer-scores}]
- Base model: 47 / 50
- Fine-tuned model: 49 / 50
\end{lstlisting}
\end{minipage}\par

微调再次带来了小幅改善。这次改善的是回答能否回应查询，可能是其中最重要的一项指标。

最后一项评估指标是答案忠实度，也称“有据性”。它检查生成的回答是否以上下文为依据。这里从上下文生成回答的是 GPT-4，向量索引只负责检索，因此可预期两个模型的表现都不错，而且差异不大。结果符合预期。

\noindent\begin{minipage}{\linewidth}
\begin{lstlisting}[caption={有据性分数},label={lst:ch07-grounding-scores}]
- Base model: 48 / 50
- Fine-tuned model: 49 / 50
\end{lstlisting}
\end{minipage}\par

表~\ref{tab:ch07-tutorial-results} 汇总结果，并列出另外两种配置，即基础版和微调版分别结合重排序。先返回前 20 份文档，再重排序并取前两份。重排序明显改善了上下文相关性，命中数接近 100\%，总 MRR 分数也略有改善。不过，重排序使答案相关性和有据性下降。原因尚不清楚，但这一结果说明，加入重排序模块时，应检验它在当前应用中的作用。

\begin{table}[!htbp]
\centering
\caption{基础版、微调版、基础版加重排序和微调版加重排序在 TruLens 三项 RAG 指标上的结果，满分为 50。}
\label{tab:ch07-tutorial-results}
\small
\setlength{\tabcolsep}{3pt}
\begin{tabular}{@{}lrrrr@{}}
\toprule
模型 & \multicolumn{2}{c}{上下文相关性} & 答案相关性 & 有据性\\
\cmidrule(lr){2-3}
 & MRR & 命中数 & &\\
\midrule
基础版 & 36.5 & 42.0 & 47 & 48\\
微调版 & 40.0 & 46.0 & 49 & 49\\
基础版加重排序 & 37.5 & 49.0 & 43 & 47\\
微调版加重排序 & 40.7 & 49.0 & 47 & 48\\
\bottomrule
\end{tabular}
\end{table}

\subsection{结论}
本练习讨论了搭建一个具备基本功能的 RAG 应用的步骤，试用了较先进的方法来改善结果，并演示怎样评估回答，判断应用是否达到预期效果。LlamaIndex 等工具可以在投入较少工作、无需大量训练模型的情况下，丰富大语言模型可用的知识。

% PDF 339--341。
\begin{chapterreferences}
\refentry{ch07-ref-001}{Vaibhav Adlakha, Parishad BehnamGhader, Xing Han Lu, Nicholas Meade, and Siva Reddy. Evaluating correctness and faithfulness of instruction-following models for question answering, 2023.}
\refentry{ch07-ref-002}{Jiawei Chen, Hongyu Lin, Xianpei Han, and Le Sun. Benchmarking large language models in retrieval-augmented generation, 2023a.}
\refentry{ch07-ref-003}{Jifan Chen, Grace Kim, Aniruddh Sriram, Greg Durrett, and Eunsol Choi. Complex claim verification with evidence retrieved in the wild, 2023b.}
\refentry{ch07-ref-004}{Shahul Es, Jithin James, Luis Espinosa-Anke, and Steven Schockaert. Ragas: Automated evaluation of retrieval augmented generation, 2023.}
\refentry{ch07-ref-005}{Zhangyin Feng, Xiaocheng Feng, Dezhi Zhao, Maojin Yang, and Bing Qin. Retrieval-generation synergy augmented large language models, 2023.}
\refentry{ch07-ref-006}{Luyu Gao, Xueguang Ma, Jimmy Lin, and Jamie Callan. Precise zero-shot dense retrieval without relevance labels. In Anna Rogers, Jordan Boyd-Graber, and Naoaki Okazaki, editors, \textit{Proceedings of the 61st Annual Meeting of the Association for} \textit{Computational Linguistics (Volume 1: Long Papers)}, pages 1762–1777, Toronto, Canada, July 2023. Association for Computational Linguistics. doi: 10.18653/v1/2023.acl-long.99. URL \url{https://aclanthology.org/2023.acl-long.99}.}
\refentry{ch07-ref-007}{Yunfan Gao et al. Retrieval-augmented generation for large language models: A survey, 2024.}
\refentry{ch07-ref-008}{Or Honovich, Leshem Choshen, Roee Aharoni, Ella Neeman, Idan Szpektor, and Omri Abend. Q\textsuperscript{2}: Evaluating factual consistency in knowledge-grounded dialogues via question generation and question answering. In Marie-Francine Moens, Xuanjing Huang, Lucia Specia, and Scott Wen-tau Yih, editors, \textit{Proceedings of the} \textit{2021 Conference on Empirical Methods in Natural Language Processing}, pages 7856–7870, Online and Punta Cana, Dominican Republic, November 2021. Association for Computational Linguistics. doi: 10.18653/v1/2021.emnlp-main.619. URL \url{https://aclanthology.org/2021.emnlp-main.619}.}
\refentry{ch07-ref-009}{Yasuto Hoshi, Daisuke Miyashita, Youyang Ng, Kento Tatsuno, Yasuhiro Morioka, Osamu Torii, and Jun Deguchi. Ralle: A framework for developing and evaluating retrieval-augmented large language models, 2023.}
\refentry{ch07-ref-010}{Chenxu Hu, Jie Fu, Chenzhuang Du, Simian Luo, Junbo Zhao, and Hang Zhao. Chatdb: Augmenting llms with databases as their symbolic memory, 2023.}
\refentry{ch07-ref-011}{Patrick Lewis, Ethan Perez, Aleksandra Piktus, Fabio Petroni, Vladimir Karpukhin, Naman Goyal, Heinrich Küttler, Mike Lewis, Wen-tau Yih, Tim Rocktäschel, et al. Retrieval-augmented generation for knowledge-intensive nlp tasks. \textit{Advances in} \textit{Neural Information Processing Systems}, 33:9459–9474, 2020.}
\refentry{ch07-ref-012}{Jerry Liu. LlamaIndex, 11 2022. URL \url{https://github.com/jerryjliu/llama_index}.}
\refentry{ch07-ref-013}{Nelson F. Liu, Kevin Lin, John Hewitt, Ashwin Paranjape, Michele Bevilacqua, Fabio Petroni, and Percy Liang. Lost in the middle: How language models use long contexts, 2023.}
\refentry{ch07-ref-014}{Xinbei Ma, Yeyun Gong, Pengcheng He, Hai Zhao, and Nan Duan. Query rewriting in retrieval-augmented large language models. In Houda Bouamor, Juan Pino, and Kalika Bali, editors, \textit{Proceedings of the 2023 Conference on Empirical Methods} \textit{in Natural Language Processing}, pages 5303–5315, Singapore, December 2023a. Association for Computational Linguistics. doi: 10.18653/v1/2023.emnlp-main.322. URL \url{https://aclanthology.org/2023.emnlp-main.322}.}
\refentry{ch07-ref-015}{Yubo Ma, Yixin Cao, YongChing Hong, and Aixin Sun. Large language model is not a good few-shot information extractor, but a good reranker for hard samples! \textit{arXiv preprint arXiv:2303.08559}, 2023b.}
\refentry{ch07-ref-016}{Yuning Mao, Pengcheng He, Xiaodong Liu, Yelong Shen, Jianfeng Gao, Jiawei Han, and Weizhu Chen. Generation-augmented retrieval for open-domain question answering. In Chengqing Zong, Fei Xia, Wenjie Li, and Roberto Navigli, editors, \textit{Proceedings of the 59th Annual Meeting of the Association for Computational Linguistics and the 11th International Joint Conference on Natural Language Processing (Volume 1: Long Papers)}, pages 4089–4100, Online, August 2021. Association for Computational Linguistics. doi: 10.18653/v1/2021.acl-long.316. URL \url{https://aclanthology.org/2021.acl-long.316}.}
\refentry{ch07-ref-017}{Niklas Muennighoff, Nouamane Tazi, Loïc Magne, and Nils Reimers. Mteb: Massive text embedding benchmark, 2023.}
\refentry{ch07-ref-018}{Colin Raffel, Noam Shazeer, Adam Roberts, Katherine Lee, Sharan Narang, Michael Matena, Yanqi Zhou, Wei Li, and Peter J. Liu. Exploring the limits of transfer learning with a unified text-to-text transformer, 2020.}
\refentry{ch07-ref-019}{Josh Reini et al. truera/trulens: Trulens eval v0.25.1, 2024. URL \url{https://zenodo.org/doi/10.5281/zenodo.4495856}.}
\refentry{ch07-ref-020}{Zhihong Shao, Yeyun Gong, Yelong Shen, Minlie Huang, Nan Duan, and Weizhu Chen. Enhancing retrieval-augmented large language models with iterative retrieval-generation synergy, 2023.}
\refentry{ch07-ref-021}{Guanzhi Wang, Yuqi Xie, Yunfan Jiang, Ajay Mandlekar, Chaowei Xiao, Yuke Zhu, Linxi Fan, and Anima Anandkumar. Voyager: An open-ended embodied agent with large language models, 2023a.}
\refentry{ch07-ref-022}{Liang Wang, Nan Yang, and Furu Wei. Query2doc: Query expansion with large language models. In Houda Bouamor, Juan Pino, and Kalika Bali, editors, \textit{Proceedings of the 2023 Conference on Empirical Methods in Natural Language Processing}, pages 9414–9423, Singapore, December 2023b. Association for Computational Linguistics. doi: 10.18653/v1/2023.emnlp-main.585. URL \url{https://aclanthology.org/2023.emnlp-main.585}.}
\refentry{ch07-ref-023}{Xiao Wang, Guangyao Chen, Guangwu Qian, Pengcheng Gao, Xiao-Yong Wei, Yaowei Wang, Yonghong Tian, and Wen Gao. Large-scale multi-modal pre-trained models: A comprehensive survey. \textit{Machine Intelligence Research}, pages 1–36, 2023c.}
\refentry{ch07-ref-024}{Xintao Wang, Qianwen Yang, Yongting Qiu, Jiaqing Liang, Qianyu He, Zhouhong Gu, Yanghua Xiao, and Wei Wang. Knowledgpt: Enhancing large language models with retrieval and storage access on knowledge bases, 2023d.}
\refentry{ch07-ref-025}{Shitao Xiao, Zheng Liu, Peitian Zhang, and Niklas Muennighoff. C-pack: Packaged resources to advance general chinese embedding, 2023.}
\refentry{ch07-ref-026}{Fangyuan Xu, Weijia Shi, and Eunsol Choi. Recomp: Improving retrieval-augmented lms with compression and selective augmentation, 2023.}
\end{chapterreferences}

